\documentclass[11pt]{article}
\usepackage[letterpaper,margin=1in,headheight=15pt]{geometry}
\usepackage[T1]{fontenc}
\usepackage{lmodern,microtype}
\usepackage{amsmath,amssymb,amsthm,mathtools}
\usepackage{booktabs,longtable,array}
\usepackage[dvipsnames]{xcolor}
\usepackage{xurl}
\usepackage[colorlinks=true,linkcolor=MidnightBlue,citecolor=MidnightBlue,urlcolor=MidnightBlue]{hyperref}
\usepackage{fancyhdr}
\usepackage{enumitem}
\usepackage{listings}
\hypersetup{pdftitle={Exact Inversion of the Partition Function: A Proof of OEIS A306631 and Sharp Rounding Bounds},pdfauthor={Research report prepared with ChatGPT},pdfsubject={Partition inversion, certified rounding, convergent corrections, and arithmetic exponential terms}}
\pagestyle{fancy}
\fancyhf{}
\fancyhead[L]{\small Exact inversion of the partition function}
\fancyhead[R]{\small A306631}
\fancyfoot[C]{\thepage}
\renewcommand{\headrulewidth}{0.3pt}
\setlength{\parskip}{0.25em}
\setlength{\emergencystretch}{2em}
\setcounter{tocdepth}{2}
\numberwithin{equation}{section}
\newtheorem{theorem}{Theorem}[section]
\newtheorem{lemma}[theorem]{Lemma}
\newtheorem{proposition}[theorem]{Proposition}
\newtheorem{corollary}[theorem]{Corollary}
\theoremstyle{definition}
\newtheorem{definition}[theorem]{Definition}
\theoremstyle{remark}
\newtheorem{remark}[theorem]{Remark}
\newcommand{\OO}{\mathcal O}
\newcommand{\N}{\mathbb N}
\newcommand{\R}{\mathbb R}
\newcommand{\C}{\mathbb C}
\newcommand{\dd}{\mathrm d}
\newcommand{\rnd}{\operatorname{rnd}}
\newcommand{\coeff}[2]{[#1^{#2}]}
\newcommand{\rising}[2]{(#1)_{#2}}
\lstset{basicstyle=\ttfamily\small,breaklines=true,columns=fullflexible,keepspaces=true,frame=single,rulecolor=\color{black!20}}

\title{\textbf{Exact Inversion of the Partition Function}\\[0.4em]
\large A proof of OEIS A306631 and sharp rounding bounds\\[0.4em]
\normalsize Convergent inverse corrections, arithmetic exponential terms,\\and exact finite verification}
\author{Research report prepared for Vladimir Reshetnikov}
\date{3 October 2026}

\begin{document}
\maketitle
\thispagestyle{empty}
\begin{abstract}
Let $p(n)$ be the ordinary partition function, and let $F$ be the increasing
inverse of $H(x)=e^{\pi\sqrt{2x/3}}/(4\sqrt3\,x)$. We prove the conjecture
recorded in OEIS A306631 that rounding $F(p(n))$ recovers $n$ for every
$n>9$. More precisely, its only failures for $n\ge2$ occur at
$2,3,4,5,7,9$, whereas taking the ceiling succeeds for every $n\ge2$.
The inverse bias $d_n=n-F(p(n))$ has the sharp, unattained infimum
$\delta=1/24+3/\pi^2$. Its maximum on $n\ge2$ occurs uniquely at $2$;
its maximum on $n\ge10$ occurs uniquely at $11$, where the bias is about
$0.4984875754$. These extrema determine every constant shift that permits
uniform exact rounding, the unique minimax shift, and a two-candidate
inversion formula requiring one comparison with a partition value.

A convergent power series describes the smooth dominant-sector bias, with
an explicit geometric remainder bound. A second convergent series inverts
the corrected dominant model. We compute twelve coefficients of each,
prove a finite hierarchy of arithmetic exponential corrections, identify
the leading parity term, and deduce fixed-order eventual alternating
finite differences. The proof combines Rademacher's classical formula
with an elementary tail estimate for $n\ge400$ and a reproducible
rational-interval certificate for the remaining indices.
\end{abstract}
\vfill
\noindent\textbf{Status.} This is an unrefereed research report. The analytic
proofs and exact finite certificate are supplied, but are not formalized
in Lean or Rocq. The literature and repository checks establish the
specific conjecture being addressed, not publication priority for every
consequence. Classical partition asymptotics and Rademacher's formula are
not claimed as new.
\newpage
{\small\setlength{\parskip}{0pt}\tableofcontents}
\newpage

\section{The problem, its status, and the main results}
\subsection{A precise OEIS question}
The OEIS entry A306631 defines the nearest-integer value of an inverse
Hardy--Ramanujan approximation, expressed using the real branch
$W_{-1}$ of the Lambert function. In the entry inspected on 3 October
2026, the assertion that this sequence takes the value $n$ at the argument
$p(n)$ for every $n>9$ is explicitly labeled a conjecture
\cite{oeis306631}. Thus the question concerns an \emph{exact integer
recovery property}, not merely asymptotic inversion.

The ordinary partition function $p(n)$, OEIS A000041 \cite{oeis000041}, counts unordered
partitions of $n$ into positive integers; $p(0)=1$. Put
\begin{equation}\label{eq:constants}
 a=\frac1{24},\qquad c=\pi\sqrt{\frac23},\qquad
 \delta=a+\frac2{c^2}=\frac1{24}+\frac3{\pi^2}.
\end{equation}
For $x>0$, define
\begin{equation}\label{eq:H}
 H(x)=\frac{e^{c\sqrt x}}{4\sqrt3\,x}.
\end{equation}
The function $H$ is strictly increasing for $x>4/c^2$. Its inverse on
that branch, restricted to $y\ge2$, is
\begin{equation}\label{eq:F}
 F(y)=\frac6{\pi^2}
 W_{-1}\!\left(-\frac{\pi}{2\sqrt2\,3^{3/4}\sqrt y}\right)^2.
\end{equation}
Write
\begin{equation}\label{eq:bias}
 d_n=n-F(p(n)),\qquad n\ge2,
 \qquad \rnd(x)=\lfloor x+\tfrac12\rfloor.
\end{equation}
The rounding convention at exact half-integers will matter for one
endpoint of a shifted-rounding theorem, but not for the original
conjecture: all its relevant comparisons are strict.

\begin{theorem}[Exact rounding and sharp bias bounds]\label{thm:main}
For every integer $n\ge2$,
\begin{equation}\label{eq:ceiling}
 0<d_n<1,\qquad \lceil F(p(n))\rceil=n.
\end{equation}
The complete nearest-integer classification is
\begin{equation}\label{eq:exceptions}
 \rnd(F(p(n)))=
 \begin{cases}
 n-1,&n\in\{2,3,4,5,7,9\},\\
 n,&\text{otherwise}.
 \end{cases}
\end{equation}
Moreover, with
\begin{equation}\label{eq:betas}
 \beta_2=2-F(2),\qquad \beta_{11}=11-F(56),
\end{equation}
one has the sharp inequalities
\begin{align}
 \delta<d_n\le\beta_2&& (n\ge2),\label{eq:globalextreme}\\
 \delta<d_n\le\beta_{11}&& (n\ge10).\label{eq:tailextreme}
\end{align}
Equality on the right holds only at $n=2$ and $n=11$, respectively.
In both ranges the infimum is $\delta$, is not attained, and is approached
as $n\to\infty$.
\end{theorem}

The conjectured threshold $10$ is therefore sharp: $9$ really is a
failure. More surprisingly, the index closest to failure \emph{after}
the threshold is $11$, not $10$. Certified enclosures give
\begin{align*}
 \delta&=0.34563021759367998099830505629584958\ldots,\\
 \beta_2&=0.79661548126519453324805686994201716\ldots,\\
 \beta_{11}&=0.49848757539660686901336470177641810\ldots.
\end{align*}
Sections~\ref{sec:tail}--\ref{sec:finite} prove the theorem. The reader
interested only in resolving A306631 may stop after that proof.

\subsection{What is established here, and what is classical}
Rademacher's convergent formula is the external analytic input; we use
the normalization in Johansson \cite{johansson}. Effective asymptotic
expansions of the partition function itself have an extensive literature;
Banerjee, Paule, Radu and Schneider \cite{banerjee} are a relevant modern
reference. None of that machinery is presented here as a new discovery.
The central deliverables of this report are the exact rounding
classification, the sharp extrema and calibration statements, and
explicitly controlled inverse formulas derived from that input.

The ProveIt repository already contains work on all-orders asymptotics,
inverse errors and integer thresholds. In particular, the inspected
A292507 report \cite{repo292507} treats the binomial transform of
partitions without ones. The present problem continues that methodology
but concerns ordinary partition values and a \emph{global exact}
rounding statement. The repository snapshot used for the source review
was
\begin{center}
\texttt{d68b65ea064084fb121df9bb431b21d086f7be81}.
\end{center}
A targeted repository search for A306631 returned no match. This is a
limited search result, not a guarantee that the subject is absent from
every repository file or from the literature. No repository theorem is
assumed in the proofs below, and no repository file has been modified.

\subsection{Why ordinary asymptotic equivalence is insufficient}
From $p(n)\sim H(n)$ one obtains only a coarse inverse approximation.
The derivative
\[
 (\log H)'(n)\sim\frac{c}{2\sqrt n}
\]
tends to zero. A relative forward error of size $n^{-1/2}$ can therefore
produce a nonzero constant inverse error. In fact the error tends to
$\delta$, rather than to zero. A proof that rounding succeeds must place
the \emph{entire} error sequence on the correct side of $1/2$; the
exceptional behavior at small odd indices cannot be discarded.

\section{The inverse branch and the exact partition expansion}
\subsection{Branch selection}
Logarithmic differentiation gives
\begin{equation}\label{eq:logderivative}
 (\log H)'(x)=\frac{c}{2\sqrt x}-\frac1x
 =\frac{c^2(t-2)}{2t^2},\qquad t=c\sqrt x.
\end{equation}
The minimum of $H$ is $c^2e^2/(16\sqrt3)<2$, so the increasing inverse
exists for all $y\ge2$. The elementary bounds $\pi^2<10$, $e<11/4$ and
$\sqrt3>5/3$ already imply that this minimum is less than $121/64<2$.
Solving $H(x)=y$ with $t=c\sqrt x>2$ gives
\[
 -\frac t2\,e^{-t/2}=-\frac{c}{4\,3^{1/4}\sqrt y}.
\]
The choice $t>2$ selects $W_{-1}$ and proves \eqref{eq:F}. The principal
Lambert branch would instead select the decreasing branch and would not
solve the OEIS question.

\subsection{The external theorem and our normalization}
For $n\ge1$ write
\begin{equation}\label{eq:mu}
 m=n-a,\qquad \mu=c\sqrt m,\qquad J(t)=t\cosh t-\sinh t.
\end{equation}
For relatively prime integers $h,k$, the Dedekind sum is
\[
 s(h,k)=\sum_{r=1}^{k-1}
 \left(\!\left(\frac rk\right)\!\right)
 \left(\!\left(\frac{hr}k\right)\!\right),
\]
where $((u))=u-\lfloor u\rfloor-1/2$ for nonintegral $u$, and $((u))=0$
for integral $u$. Define
\begin{equation}\label{eq:Ak}
 A_k(n)=\sum_{\substack{0\le h<k\\(h,k)=1}}
 \exp\!\left(\pi i\left(s(h,k)-\frac{2nh}{k}\right)\right).
\end{equation}
For $k=1$ the convention yields $A_1(n)=1$; also
$A_2(n)=(-1)^n$. Termwise absolute values give $|A_k(n)|\le k$.
Each $A_k(n)$ is periodic in $n$ modulo $k$.

Rademacher's formula, in the equivalent normalization of
\cite[Section~1]{johansson}, is
\begin{equation}\label{eq:rademacher}
 p(n)=\frac1{2\pi\sqrt2\,m^{3/2}}
 \sum_{k\ge1}A_k(n)\sqrt k\,J(\mu/k).
\end{equation}
The classical identity \eqref{eq:rademacher} is quoted, not reproved.
The estimates, inversions and finite verifications that follow are
proved within this report.

\section{An elementary effective Rademacher tail}\label{sec:tail}
\begin{lemma}[Absolute tail bound]\label{lem:tail}
For every $K\ge1$ and $n\ge1$, if $T_k(n)$ is the $k$th term of
\eqref{eq:rademacher}, then
\begin{equation}\label{eq:tail}
 \left|p(n)-\sum_{k=1}^K T_k(n)\right|
 \le\frac{2\pi^2}{9\sqrt{3K}}\,e^{\mu/(K+1)}.
\end{equation}
\end{lemma}
\begin{proof}
For $t\ge0$,
\[
 J(t)=\int_0^t u\sinh u\,\dd u
 \le\int_0^t u^2e^u\,\dd u\le\frac{t^3e^t}{3}.
\]
Consequently
\[
 |T_k(n)|\le\frac{c^3}{6\pi\sqrt2}
 k^{-3/2}e^{\mu/k}
 =\frac{\pi^2}{9\sqrt3}k^{-3/2}e^{\mu/k}.
\]
This also proves absolute convergence for fixed $n$. For $k>K$ use
$e^{\mu/k}\le e^{\mu/(K+1)}$ and
$\sum_{k>K}k^{-3/2}\le\int_K^\infty x^{-3/2}\dd x=2/\sqrt K$.
\end{proof}

The bound is deliberately elementary rather than optimal. Its purpose is
to turn an infinite-index claim into a finite rational certificate, not
to compete with optimized implementations of Rademacher's formula.

Define the positive dominant sector, extended to real $x>a+c^{-2}$, by
\begin{equation}\label{eq:M}
 M(x)=\frac{e^{c\sqrt{x-a}}}{4\sqrt3\,(x-a)}
 \left(1-\frac1{c\sqrt{x-a}}\right).
\end{equation}
Splitting the $k=1$ summand into its positive and negative exponentials
shows that
\[
 T_1(n)=M(n)+\frac{e^{-\mu}}{4\sqrt3\,m}(1+\mu^{-1}).
\]
\begin{lemma}[Relative error of the dominant sector]\label{lem:relative}
For $\mu\ge50$, write $p(n)=M(n)(1+\epsilon_n)$. Then
\begin{equation}\label{eq:epsilon}
 |\epsilon_n|\le\frac32\mu^2e^{-\mu/2}.
\end{equation}
\end{lemma}
\begin{proof}
The $K=1$ case of \eqref{eq:tail} and the displayed negative exponential
term imply
\begin{equation}\label{eq:epsilonprecise}
 |\epsilon_n|\le
 \frac{1+1/\mu}{1-1/\mu}e^{-2\mu}
 +\frac43\frac{\mu^2}{1-1/\mu}e^{-\mu/2}.
\end{equation}
For $\mu\ge50$, the coefficient of $\mu^2e^{-\mu/2}$ in the second
term is at most $200/147$. The first term divided by that same scale
is at most $51/122500$, even after discarding its additional decaying
exponential. The sum of those rational bounds is less than $3/2$.
\end{proof}

\section{A convergent model for the inverse bias}\label{sec:smooth}
The dominant sector can be inverted without making a merely formal
asymptotic ansatz. Its bias is an analytic function of $1/\mu$.

\begin{lemma}[Analytic fixed point]\label{lem:v}
There is a unique analytic solution with $|v(z)-1|\le1/4$ on
$|z|\le1/20$ of
\begin{equation}\label{eq:vfixed}
 v(z)=-\frac{\log(1-z)}z
       -\frac2z\log(1-z^2v(z)),\qquad v(0)=1.
\end{equation}
The apparent singularities at $z=0$ are removable. All Taylor
coefficients of $v$ are nonnegative. In particular, for $0\le z\le1/20$,
\begin{equation}\label{eq:vrealbounds}
 1+\frac52 z\le v(z)\le\frac54.
\end{equation}
\end{lemma}
\begin{proof}
Use the analytic logarithms determined by their series at $1$. Put
$r=1/20$ and $Q=5/4$. On $|z|\le r$, $|v-1|\le1/4$, the right-hand
side of \eqref{eq:vfixed}, minus $1$, has modulus at most
\[
 \frac{r}{2(1-r)}+\frac{2rQ}{1-r^2Q}<\frac14.
\]
Its derivative with respect to $v$ has modulus at most
\[
 \frac{2r}{1-r^2Q}<\frac14.
\]
Uniform fixed-point iteration starting at $v=1$ is thus a contraction
on that disk, and converges uniformly to an analytic solution. The same
bounds hold in a slightly larger $z$-disk, so analyticity on a
neighborhood of the closed disk is available for Cauchy estimates.

Expanding \eqref{eq:vfixed} gives
\begin{equation}\label{eq:vrec}
 v_0=1,\qquad
 v_j=\frac1{j+1}
 +2\sum_{k=1}^{\lfloor(j+1)/2\rfloor}
 \frac1k\coeff{z}{j-2k+1}v(z)^k\quad(j\ge1).
\end{equation}
Only earlier coefficients occur on the right, and they occur with
nonnegative coefficients. Induction proves $v_j\ge0$, while
$v_1=5/2$. The upper real bound is the fixed-point disk bound.
\end{proof}

\begin{definition}\label{def:B}
For $|z|\le1/20$, set
\begin{equation}\label{eq:B}
 B(z)=a+\frac{2v(z)-z^2v(z)^2}{c^2}.
\end{equation}
\end{definition}
\begin{proposition}[Exact dominant-sector identity]\label{prop:Bidentity}
Whenever $\mu=c\sqrt{n-a}\ge50$,
\begin{equation}\label{eq:Bidentity}
 n-F(M(n))=B(1/\mu).
\end{equation}
The same identity holds for real $n$ in this range.
\end{proposition}
\begin{proof}
Set $z=1/\mu$ and $\nu=\mu-zv(z)$. Equation~\eqref{eq:vfixed} is
exactly
\[
 -zv-2\log(1-z^2v)=\log(1-z).
\]
Hence
\[
 \frac{e^\nu}{\nu^2}
 =\frac{e^\mu}{\mu^2}(1-\mu^{-1}).
\]
Since $\nu>2$, the increasing inverse branch gives
$F(M(n))=\nu^2/c^2$. Subtracting this quantity from
$n=a+\mu^2/c^2$ proves the assertion.
\end{proof}

\begin{remark}
The exact identity concerns $M(n)$, not $p(n)$. The distinction is
essential: $B(1/\mu)$ is a convergent analytic main sector, whereas
$d_n$ also contains arithmetic exponentially small terms. These terms
will be made explicit in Section~\ref{sec:arithmetic}.
\end{remark}

\section{Uniform inequalities for every large index}\label{sec:effective}
\begin{theorem}[Effective inverse estimate]\label{thm:effective}
For every integer $n\ge400$, with $\mu=c\sqrt{n-a}$,
\begin{equation}\label{eq:inverseerror}
 |d_n-B(1/\mu)|\le
 \frac{9\mu^3}{c^2}e^{-\mu/2}<\frac1{c^2\mu}.
\end{equation}
In particular,
\begin{equation}\label{eq:largebracket}
 \delta+\frac3{c^2\mu}<d_n<\frac{43}{100},
 \qquad \lim_{n\to\infty}d_n=\delta.
\end{equation}
\end{theorem}
\begin{proof}
The elementary bounds $25/8<\pi<22/7$ imply $625/96<c^2<7$.
Thus $n\ge400$ gives
$\mu^2>(625/96)(400-1/24)>2500$, so $\mu\ge50$.
Let $z=1/\mu$. From \eqref{eq:vrealbounds},
\begin{equation}\label{eq:Bbracket}
 B(z)-\delta\ge\frac{5z-(25/16)z^2}{c^2},\qquad
 0<B(z)\le a+\frac{5}{2c^2}<\frac{43}{100}.
\end{equation}
Write $x_0=F(M(n))=n-B(z)$. We first locate $F(p(n))$ in
$[n-1,n+1]$ rather than presupposing that location.

For $x\in[n-1,n+1]$ put $t=c\sqrt x$. Since $c^2<7$ and $\mu\ge50$,
we have $\mu-1<t<\mu+1$. The function $(t-2)/t^2$ decreases for
$t>4$. Equation~\eqref{eq:logderivative} therefore gives
\begin{equation}\label{eq:derivativebound}
 (\log H)'(x)\ge\frac{c^2(\mu-1)}{2(\mu+1)^2}
 \ge\frac{c^2}{3\mu}.
\end{equation}
The last inequality is equivalent to
$\mu^2-7\mu-2\ge0$.

The function $\mu^4e^{-\mu/2}$ decreases for $\mu\ge50$, and
\begin{equation}\label{eq:expaux}
 e^{25}>9\cdot50^4.
\end{equation}
In fact $e^{25}>25^{11}/11!>9\cdot50^4$, a comparison of integers
after multiplication by $11!$. The accompanying exact certificate also
checks the inequality. Consequently $9\mu^4e^{-\mu/2}<1$ throughout the range.
Lemma~\ref{lem:relative} now gives $|\epsilon_n|<1/2$ and
\begin{equation}\label{eq:logepsilon}
 |\log(1+\epsilon_n)|\le2|\epsilon_n|
 \le3\mu^2e^{-\mu/2}<\frac1{3\mu^2}.
\end{equation}

The point $x_0$ is more than $1/2$ from either endpoint of $[n-1,n+1]$,
by \eqref{eq:Bbracket}. Integrating \eqref{eq:derivativebound} from
$x_0$ to those endpoints shows that changing $\log M(n)$ by the amount
in \eqref{eq:logepsilon} cannot cross either endpoint value of $\log H$;
indeed $1/(3\mu^2)<c^2/(6\mu)$. Thus $F(p(n))$ lies inside that
interval. Applying the mean value theorem there proves
\[
 |F(p(n))-F(M(n))|
 \le\frac{3\mu}{c^2}|\log(1+\epsilon_n)|
 \le\frac{9\mu^3}{c^2}e^{-\mu/2}.
\]
Together with \eqref{eq:expaux}, this is \eqref{eq:inverseerror}.

Subtracting $z/c^2$ from the lower bound in \eqref{eq:Bbracket} gives
\[
 d_n-\delta>\frac{4z-(25/16)z^2}{c^2}
 >\frac{3z}{c^2}.
\]
For the upper bound use $z\le1/50$ to obtain
\[
 d_n<a+\frac{5/2+1/50}{c^2}
 =a+\frac{63}{25c^2}<\frac{43}{100}.
\]
For the last inequality, $c^2>625/96$ gives
$a+63/(25c^2)<160777/375000<43/100$.
The same elementary bound gives
$\delta<5233/15000<7/20$. These are also checked in the certificate. Finally, $B(0)=\delta$
and the error in \eqref{eq:inverseerror} tends to zero.
\end{proof}

\section{The exact finite certificate and completion of the proof}\label{sec:finite}
\subsection{What must be checked}
Every comparison below takes place at $x\ge1$, where $H$ is strictly
increasing. Thus
\begin{equation}\label{eq:comparisonrule}
 d_n>s\ \Longleftrightarrow\ p(n)<H(n-s),\qquad
 d_n<s\ \Longleftrightarrow\ p(n)>H(n-s).
\end{equation}
This avoids evaluating a Lambert function in any proof check.

\begin{lemma}[Certified finite inequalities]\label{lem:finite}
The following strict inequalities hold:
\begin{align}
 \frac7{20}<d_n<1&& (2\le n\le399),\label{eq:finiteall}\\
 d_n<\frac{49}{100}&& (10\le n\le399,\ n\ne11),\label{eq:finite49}\\
 \frac{498}{1000}<d_{11}<\frac{499}{1000},\qquad
 &\frac{796}{1000}<d_2<\frac{797}{1000},\label{eq:finitespecial}\\
 d_n<\frac{79}{100}&& (3\le n\le9).\label{eq:finite79}
\end{align}
For $2\le n\le9$, the inequality $d_n>1/2$ holds exactly at
$n=2,3,4,5,7,9$; at $6,8$ one has $d_n<1/2$.
\end{lemma}
\begin{proof}[Proof by exact finite computation]
The file \texttt{code/certify.py} implements \eqref{eq:comparisonrule}
using only integer arithmetic, integer square roots and rational
endpoints. It completed all 1,204 asserted comparisons, with every
interval separated strictly from the comparison value. The independent
recomputation of all partition values also agreed. The construction of
the intervals, including their rigorous transcendental remainders, is
given below so that the finite computation can be audited without
trusting a floating-point library. Its machine-readable outcome is
\texttt{results/certificate.json}.
\end{proof}

\subsection{Why the intervals are rigorous}
All stored interval endpoints are integers divided by
$S=10^{70}$. Addition is exact at this scale; multiplication and division
round lower endpoints downward and upper endpoints upward. Square roots
use integer square roots of scaled nonnegative integers, with an extra
unit on an upper endpoint unless the square is exact. These operations
enclose the corresponding real operations.

For $\pi$, the program uses Machin's identity
\[
 \pi=16\arctan(1/5)-4\arctan(1/239).
\]
Each arctangent is enclosed between its alternating-series partial sum
and that sum plus the next term. The program uses 60 terms and 15 terms,
respectively, before rounding the rational endpoints outward.
The identity follows by the tangent addition formula, with the relevant
angles in the first quadrant.

For $0\le t\le64$, it computes $\exp(t/64)$ through degree 80 and then
squares the interval six times. On $0\le u\le1$ the omitted tail is
bounded by
\[
 \sum_{j=81}^\infty\frac{u^j}{j!}
 \le\frac1{81!}\sum_{r\ge0}82^{-r}
 =\frac{82}{81\cdot81!}<\frac2{81!}<\frac1S.
\]
Adding one scaled unit to the upper endpoint is therefore sufficient
before the six squarings. Every exponential argument required by the
finite comparisons is below 64. The computation of $H$ then uses only
these operations on rational arguments.

Partition numbers are first computed by Euler's recurrence
\begin{equation}\label{eq:pentagonal}
 p(n)=\sum_{j\ge1}(-1)^{j-1}
 \left(p\!\left(n-\frac{j(3j-1)}2\right)
      +p\!\left(n-\frac{j(3j+1)}2\right)\right),
\end{equation}
with $p(0)=1$ and $p(t)=0$ for $t<0$. As an independent check, a second
program loop multiplies the truncated factors $(1-q^j)^{-1}$ through
$q^{399}$. Equivalently, starting with $b_0=1$ and $b_n=0$ for $n>0$,
for each part size $j=1,\ldots,399$ it performs $b_n\leftarrow b_n+b_{n-j}$
in increasing order of $n\ge j$. This enumerates the possible
multiplicities of each part size. All 400 coefficients agree with
\eqref{eq:pentagonal}. Thus even the pentagonal identity is not needed
as an unchecked source of the finite data.

The asserted $H$ comparisons are grouped as follows.
\begin{center}
\begin{tabular}{@{}p{0.73\linewidth}r@{}}
\toprule
Claim checked & Comparisons\\
\midrule
$7/20<d_n<1$ for $2\le n\le399$ & 796\\
$d_n<49/100$ for $10\le n\le399$, $n\ne11$ & 389\\
Two-sided brackets at $n=2$ and $n=11$ & 4\\
$d_n<79/100$ for $3\le n\le9$ & 7\\
Comparison with $1/2$ for $2\le n\le9$ & 8\\
\midrule
Total asserted $H$ comparisons & 1,204\\
\bottomrule
\end{tabular}
\end{center}
The smallest certified absolute separation among these comparisons
exceeds $0.0002600145755$; it occurs in the check $d_2<0.797$.
Additional exact checks establish $\delta<7/20$,
\eqref{eq:expaux}, and the upper bound at the end of
Theorem~\ref{thm:effective}. Rational bisection of $H$, again without
Lambert $W$, supplies the displayed high-precision enclosures for
$\beta_2$ and $\beta_{11}$.

\subsection{Completion of the main theorem}
\begin{proof}[Proof of Theorem~\ref{thm:main}]
For $n\ge400$, Theorem~\ref{thm:effective} gives
$\delta<d_n<0.43$. For $2\le n\le399$, Lemma~\ref{lem:finite} gives
$0.35<d_n<1$, and $\delta<0.35$. This proves $0<d_n<1$ and hence the
ceiling identity.

For $n\ge10$, the finite bound is below $0.49$ except at $11$, where
it is below $0.499$; the infinite tail is below $0.43$. Thus $d_n<1/2$
for all $n\ge10$. The last assertion of Lemma~\ref{lem:finite}
classifies the indices from $2$ to $9$. Since $0<d_n<1$, their rounded
values are exactly those in \eqref{eq:exceptions}.

The value $d_{11}>0.498$ is larger than every other $d_n$ on $n\ge10$:
those with $n<400$ are below $0.49$, and the rest are below $0.43$.
Similarly, $d_2>0.796$ exceeds the bounds $0.79$ for $3\le n\le9$
and $0.499$ for $n\ge10$. Thus both maxima are unique.
All $d_n$ exceed $\delta$, while Theorem~\ref{thm:effective} gives
$d_n\to\delta$. This proves the sharp, unattained infima.
\end{proof}

\begin{remark}[Scope of the certificate]
This is an exact-arithmetic computer-assisted proof of a finite lemma,
not a numerical experiment at selected large indices. The infinite
range has a separate analytic proof. Nevertheless, executing Python
and checking its source is not the same as a small-kernel formal proof;
no Lean or Rocq certification is claimed. Reproduction should use
ordinary Python execution, not the \texttt{-O} flag, which removes
assertions.
\end{remark}

\section{Optimal constant shifts and robust rounding}\label{sec:calibration}
The sharp extrema provide more information than the truth of the
original conjecture.

\begin{theorem}[All successful shifts and the minimax shift]\label{thm:shifts}
Let $S$ be either $\{n\in\N:n\ge2\}$ or $\{n\in\N:n\ge10\}$, and let
$\beta$ be $\beta_2$ or $\beta_{11}$, respectively. For a real constant
$s$, the identity
\begin{equation}\label{eq:shiftedround}
 \rnd(F(p(n))+s)=n\qquad(n\in S)
\end{equation}
holds if and only if
\begin{equation}\label{eq:shiftinterval}
 \beta-\frac12\le s\le\delta+\frac12.
\end{equation}
The unique constant shift minimizing the worst absolute real error is
\begin{equation}\label{eq:minimax}
 s_* =\frac{\beta+\delta}{2},\qquad
 \inf_{s\in\R}\sup_{n\in S}|F(p(n))+s-n|
 =\frac{\beta-\delta}{2}.
\end{equation}
\end{theorem}
\begin{proof}
By the rounding convention, \eqref{eq:shiftedround} is equivalent to
$-1/2\le s-d_n<1/2$ for every $n\in S$. The left inequalities require
$s\ge\sup d_n-1/2=\beta-1/2$. The right inequalities allow precisely
$s\le\delta+1/2$: equality is allowed because $d_n>\delta$ for every
finite $n$, while any larger $s$ fails along a sequence tending to the
infimum. This proves \eqref{eq:shiftinterval}.

Since the range has supremum $\beta$ and infimum $\delta$,
\[
 \sup_{n\in S}|s-d_n|=\max\{|s-\delta|,|s-\beta|\}.
\]
The unique minimizer is the midpoint, with half the interval length as
its value.
\end{proof}

For $n\ge10$, exact interval bisection gives
\begin{align*}
 s_*&=0.42205889649514342500583487903613384\ldots,\\
 E_*&=0.07642867890146344400752982274028426\ldots.
\end{align*}
Thus a fixed shift reduces the worst real error on the conjectured
range from almost $1/2$ to less than $0.077$. The set of all successful
shifts is approximately
\[
 [-0.00151242460339313\ldots,\ 0.84563021759367998\ldots]
 \qquad(n\ge10).
\]
For all $n\ge2$, the lower endpoint instead becomes
$0.29661548126519453\ldots$. In particular $s=1/2$ works for every
$n\ge2$, which is another expression of the ceiling theorem.

The smallest symmetric numerical-error margin of the unshifted
formula on $n\ge10$ is
\begin{equation}\label{eq:roundmargin}
 \rho=\frac12-\beta_{11}
 =0.00151242460339313098663529822358189\ldots.
\end{equation}
An evaluation error of absolute value less than $\rho$ cannot change
any of those rounded answers. No larger uniform symmetric margin can
be valid, because the lower rounding boundary at $n=11$ is exactly
$\rho$ away. At the boundary itself the result depends on the tie
convention; our lower endpoint in \eqref{eq:shiftinterval} uses
rounding half-integers upward.

\section{Inverting the entire partition staircase}\label{sec:staircase}
Exact inversion at partition values is different from inversion at an
arbitrary real target. Define
\begin{equation}\label{eq:N}
 N(y)=\min\{n\ge0:p(n)\ge y\},\qquad y\ge42.
\end{equation}
Here the minimum exists because $p(n)$ is unbounded, and the relevant
partition values are strictly increasing. Adding a part $1$ injects
partitions of $n$ into those of $n+1$, while the one-part partition of
$n+1$ is missing from that image for $n\ge1$.

\begin{theorem}[Two candidates and one partition comparison]\label{thm:staircase}
For $y\ge42$, put
\begin{equation}\label{eq:k}
 k(y)=\lfloor F(y)+\delta\rfloor+1.
\end{equation}
Then
\begin{equation}\label{eq:onecomparison}
 N(y)=
 \begin{cases}
 k(y),&p(k(y))\ge y,\\
 k(y)+1,&p(k(y))<y.
 \end{cases}
\end{equation}
\end{theorem}
\begin{proof}
Suppose first that $y>42$, and set $n=N(y)\ge11$.
The inequalities $p(n-1)<y\le p(n)$ imply
\[
 n-1-d_{n-1}<F(y)\le n-d_n.
\]
Using $\delta<d_j\le\beta_{11}$ for $j\ge10$ gives
\begin{equation}\label{eq:staircaseinterval}
 F(y)+\delta<n<F(y)+1+\beta_{11}.
\end{equation}
The interval length is $1+\beta_{11}-\delta<2$. Its smallest possible
integer is $k=\lfloor F(y)+\delta\rfloor+1$, so $n\in\{k,k+1\}$.
Monotonicity of $p$ distinguishes these two possibilities by the one
comparison in \eqref{eq:onecomparison}. For $y=42$, the finite bounds on
$d_{10}$ give $9<F(42)+\delta<10$, so $k=10=N(42)$ directly.
\end{proof}

This counts evaluations of the \emph{partition function}, assuming the
integer $k$ has been computed correctly. It is not a bit-complexity
claim about Lambert $W$, nor a claim that a fixed floating-point
precision always determines a floor for arbitrary real input.
Certified intervals for $F(y)+\delta$ must resolve the floor, and an
input exactly on a floor boundary needs an exact representation or an
appropriate equality convention. The numerical experiment shipped here
checks all 4,959 integer targets $42\le y\le5000$; the theorem itself
covers all real $y\ge42$.

\section{All-orders bias coefficients and an explicit remainder}\label{sec:coefficients}
The fixed-point equation gives an exact rational algorithm for every
coefficient of the analytic bias, not only its first few terms.

\begin{theorem}[Convergent bias series]\label{thm:biasseries}
On $|z|<1/20$,
\begin{equation}\label{eq:Bseries}
 B(z)=\delta+\frac1{c^2}\sum_{j\ge1}w_jz^j,
\end{equation}
where the $v_j$ are given by \eqref{eq:vrec} and
\begin{equation}\label{eq:wrec}
 w_j=2v_j-\coeff{z}{j-2}v(z)^2\qquad(j\ge1),
\end{equation}
with a negative-index coefficient interpreted as zero.
For every $J\ge0$ and real $\mu>20$,
\begin{equation}\label{eq:cauchyB}
 \left|B(1/\mu)-\delta-\frac1{c^2}
             \sum_{j=1}^Jw_j\mu^{-j}\right|
 \le\frac{129}{256c^2}
       \frac{(20/\mu)^{J+1}}{1-20/\mu}.
\end{equation}
Consequently, for $n\ge400$, the difference between $d_n$ and the same
truncated series is at most the right-hand side of
\eqref{eq:cauchyB} plus $9\mu^3e^{-\mu/2}/c^2$.
\end{theorem}
\begin{proof}
The series follows from Lemma~\ref{lem:v} and \eqref{eq:B}. On
$|z|=1/20$, that lemma gives
\[
 |c^2(B(z)-\delta)|
 =|2(v-1)-z^2v^2|
 \le\frac12+\frac1{400}\frac{25}{16}
 =\frac{129}{256}.
\]
Cauchy's coefficient estimate gives
$|w_j|\le(129/256)20^j$. Summing the resulting geometric tail proves
\eqref{eq:cauchyB}. The assertion for $d_n$ follows from
Theorem~\ref{thm:effective}.
\end{proof}

The first terms are
\begin{equation}\label{eq:forwardterms}
 d_n=\delta+\frac5{c^2\mu}
 +\frac{29}{3c^2\mu^2}+\frac{113}{6c^2\mu^3}
 +\frac{823}{20c^2\mu^4}+\OO(\mu^{-5}).
\end{equation}
More precisely, the complete convergent analytic series is accompanied
by the exponentially small remainder in \eqref{eq:inverseerror}.
In powers of $n^{-1/2}$ the first terms become
\begin{align}\label{eq:npowers}
 d_n={}&\delta+\frac5{c^3}n^{-1/2}
 +\frac{29}{3c^4}n^{-1}\notag\\
 &+\left(\frac{5a}{2c^3}+\frac{113}{6c^5}\right)n^{-3/2}
 +\OO(n^{-2}).
\end{align}
The shift $a=1/24$ must not be dropped when converting higher terms.

The finite algorithm in \texttt{code/series.py} uses exact rational
arithmetic. It is triangular: the coefficient $v_j$ in
\eqref{eq:vrec} depends only on $v_0,\ldots,v_{j-1}$, so any prescribed
finite order requires only finitely many additions and multiplications.
Positivity of $v_j$ does \emph{not}, without a further argument,
establish positivity of all $w_j$, because \eqref{eq:wrec} includes a
subtraction. The twelve positive values computed below are not used to
assert that stronger statement.

\section{A convergent corrected inverse at an arbitrary target}\label{sec:correctedinverse}
The preceding series describes the error at a known index $n$. For
inversion from a target $y$, one wants a correction expressed in terms
of $F(y)$ instead. These are different series.

For $\mu=c\sqrt{x-a}>1$, differentiation of \eqref{eq:M} yields
\begin{equation}\label{eq:Mderivative}
 (\log M)'(x)
 =\frac{c^2(\mu^2-3\mu+3)}{2\mu^2(\mu-1)}>0.
\end{equation}
The numerator is $(\mu-3/2)^2+3/4$. Since $M$ increases from $0$ to
$+\infty$ on $x>a+c^{-2}$, it has an inverse $G$ on $y>0$.

\begin{theorem}[Convergent correction to the Lambert inverse]\label{thm:q}
There is a unique analytic function $q$ satisfying $q(0)=1$,
$|q(w)-1|\le1/4$ on $|w|\le1/20$, and
\begin{equation}\label{eq:qeq}
 wq-3\log(1+w^2q)+\log(1-w+w^2q)=0.
\end{equation}
For $u=F(y)$, $t=c\sqrt u>20$, and $w=1/t$,
\begin{equation}\label{eq:Gexact}
 G(y)=u+a+\frac{2q(w)+w^2q(w)^2}{c^2}
 =u+\delta+\frac1{c^2}\sum_{j\ge1}r_jt^{-j}.
\end{equation}
The coefficients are rational and are determined successively by
\begin{align}
 q_0&=1,\notag\\
 q_j&=\coeff{w}{j+1}
 \{3\log(1+w^2q(w))-\log(1-w+w^2q(w))\},\label{eq:qrec}\\
 r_j&=2q_j+\coeff{w}{j-2}q(w)^2.\label{eq:rrec}
\end{align}
For every $J\ge0$, truncation after $r_Jt^{-J}$ has absolute error
at most
\begin{equation}\label{eq:cauchyG}
 \frac{129}{256c^2}\frac{(20/t)^{J+1}}{1-20/t}.
\end{equation}
\end{theorem}
\begin{proof}
Divide \eqref{eq:qeq} by $w$ and rearrange as the fixed-point equation
\begin{equation}\label{eq:qmap}
 q=\frac2w\log(1+w^2q)
 -\frac1w\log\!\left(1-\frac{w}{1+w^2q}\right).
\end{equation}
As before, the apparent singularities at zero are removable. For
$|w|\le r=1/20$, $|q-1|\le1/4$ and $Q=5/4$, the right-hand side minus
$1$ has modulus at most
\[
 \frac{2rQ}{1-r^2Q}
 +\frac{r^2Q}{1-r^2Q}
 +\frac{r}{2(1-r^2Q)(1-r^2Q-r)}<\frac14.
\]
Here the last two terms bound the second logarithm in
\eqref{eq:qmap} minus $1$. Differentiating the original rearrangement
of \eqref{eq:qeq} bounds the $q$-derivative by
\[
 \frac{3r}{1-r^2Q}+\frac{r}{1-r-r^2Q}<\frac14.
\]
Uniform analytic contraction proves existence and uniqueness.
Coefficient comparison gives \eqref{eq:qrec}; its right side depends
only on earlier $q$ coefficients.

For real $w=1/t>0$, set $\mu=t+wq(w)$ and $x=a+\mu^2/c^2$.
The relation $M(x)=H(u)$, with $u=t^2/c^2$, is equivalent to
\[
 (\mu-t)-2\log(\mu/t)+\log(1-\mu^{-1})=0.
\]
Substitute $\mu/t=1+w^2q$ to obtain \eqref{eq:qeq}. The uniqueness of
$G$ then gives \eqref{eq:Gexact}. Finally,
\[
 |2(q-1)+w^2q^2|\le\frac{129}{256}
\]
on $|w|=1/20$, so the same Cauchy argument as before proves
\eqref{eq:cauchyG}.
\end{proof}

The first correction terms are
\begin{equation}\label{eq:Gfirst}
 G(y)=F(y)+\delta+\frac5{c^2t}
 +\frac{29}{3c^2t^2}+\frac{83}{6c^2t^3}
 +\frac{559}{60c^2t^4}+\OO(t^{-5}).
\end{equation}
The third coefficient is $83/6$, whereas the bias coefficient in
\eqref{eq:forwardterms} is $113/6$. Interchanging them amounts to
replacing the unknown index by its uncorrected inverse inside a
higher-order term, and produces an incorrect expansion.

\begin{table}[tbp]
\centering
\caption{Exact coefficients of the two convergent inverse constructions.
Here $w_j$ belongs to $B(1/\mu)$ and $r_j$ to $G(y)-F(y)$.}\label{tab:coeffs}
\medskip
\begin{tabular}{@{}r rr@{}}
\toprule
$j$ & $w_j$ & $r_j$\\
\midrule
1 & $5$ & $5$\\
2 & $29/3$ & $29/3$\\
3 & $113/6$ & $83/6$\\
4 & $823/20$ & $559/60$\\
5 & $1517/15$ & $-611/30$\\
6 & $338257/1260$ & $-112459/1260$\\
7 & $938657/1260$ & $-101329/630$\\
8 & $10713193/5040$ & $-228677/5040$\\
9 & $15593947/2520$ & $1709167/2520$\\
10 & $317064037/17325$ & $325098619/138600$\\
11 & $19540325693/356400$ & $329128673/89100$\\
12 & $1076961671111/6486480$ & $-9006152791/6486480$\\
\bottomrule
\end{tabular}
\end{table}

At actual partition targets, the corrected model is exponentially
accurate rather than exact:
\begin{equation}\label{eq:Gactual}
 G(p(n))-n=\OO(\mu^3e^{-\mu/2}).
\end{equation}
To see this, combine \eqref{eq:epsilon} with the positive logarithmic
derivative \eqref{eq:Mderivative}, whose reciprocal is $\OO(\mu)$
on a fixed-width neighborhood of $n$. The same bracketing argument as
in Theorem~\ref{thm:effective} applies for sufficiently large $n$.
The bound in \eqref{eq:Gactual} tends to zero, so rounding $G(p(n))$
eventually recovers $n$. No sharp starting index for that separate
corrected-model rule is asserted here.

\section{Arithmetic exponential corrections}\label{sec:arithmetic}
The coefficients $A_k(n)$ create periodic fluctuations that no fixed
power series in $n^{-1/2}$ can capture. We now isolate them at the
inverse level with a controlled remainder.

\begin{theorem}[Any fixed finite set of arithmetic sectors]\label{thm:sectors}
Fix an integer $K\ge2$. Put $z=1/\mu$ and
$\nu=\mu-v(z)/\mu$. As $n\to\infty$ through integers,
\begin{align}
 d_n={}&B(1/\mu)
 -\frac{2\nu^2}{c^2(\nu-2)}
 \sum_{k=2}^K\frac{A_k(n)}{\sqrt k}
   \frac{1-k/\mu}{1-1/\mu}
   e^{-(1-1/k)\mu}\notag\\
 &+\OO_K\!\left(\mu^3
       e^{-(1-1/(K+1))\mu}\right).
 \label{eq:sectors}
\end{align}
The implied constant can depend on the fixed integer $K$, but not on
$n$ or its residue class.
\end{theorem}
\begin{proof}
Splitting $J(\mu/k)$ into exponentials, the positive-exponential part
of the $k$th Rademacher summand is
\[
 P_{k,+}(n)=\frac{A_k(n)}{4\sqrt3\,m\sqrt k}
 e^{\mu/k}(1-k/\mu).
\]
Its ratio to $M(n)$ is
\begin{equation}\label{eq:sectorratio}
 E_k(n)=\frac{A_k(n)}{\sqrt k}
 \frac{1-k/\mu}{1-1/\mu}e^{-(1-1/k)\mu}.
\end{equation}
For fixed $K$, the negative-exponential parts of the first $K$ terms,
after division by $M(n)$, are
$\OO_K(e^{-(1+1/K)\mu})$. Lemma~\ref{lem:tail} bounds the remaining
exact Rademacher tail divided by $M(n)$ by
$\OO_K(\mu^2e^{-(1-1/(K+1))\mu})$. Therefore
\begin{equation}\label{eq:epssectors}
 \epsilon_n=\sum_{k=2}^K E_k(n)
 +\OO_K(\mu^2e^{-(1-1/(K+1))\mu}).
\end{equation}
In particular, taking $K=2$ first shows
$\epsilon_n=\OO(e^{-\mu/2})$: the polynomial factor in the next
exponential is absorbed by the gap between its exponent and $1/2$.

Let $\Phi(\ell)=F(e^\ell)$. At $\ell=\log M(n)$,
\begin{equation}\label{eq:PhiDerivative}
 \Phi'(\ell)=\frac{2\nu^2}{c^2(\nu-2)}=\OO(\mu).
\end{equation}
On a bounded neighborhood of this argument, $\Phi''(\ell)=\OO(1)$.
Indeed $\ell=\nu-2\log\nu+\log(c^2/(4\sqrt3))$,
so $\dd\nu/\dd\ell=\nu/(\nu-2)$, and direct differentiation of
\eqref{eq:PhiDerivative} is bounded for large $\nu$.
Taylor expansion in $h=\log(1+\epsilon_n)$ gives
\[
 F(p(n))-F(M(n))
 =\frac{2\nu^2}{c^2(\nu-2)}\epsilon_n
 +\OO(\mu\epsilon_n^2).
\]
The nonlinear error is $\OO(\mu e^{-\mu})$, smaller than the stated
remainder for every fixed $K$. Substitution of \eqref{eq:epssectors}
and the identity $d_n=B-(F(p(n))-F(M(n)))$ proves the result.
\end{proof}

\begin{corollary}[Leading parity effect]\label{cor:parity}
Define the analytic amplitude near $z=0$ by
\begin{equation}\label{eq:amplitude}
 \mathcal A(z)=
 \frac{(1-z^2v(z))^2(1-2z)}
 {(1-2z-z^2v(z))(1-z)},\qquad\mathcal A(0)=1.
\end{equation}
Then
\begin{equation}\label{eq:parity}
 d_n=B(1/\mu)
 -(-1)^n\frac{\sqrt2\,\mu}{c^2}
 \mathcal A(1/\mu)e^{-\mu/2}
 +\OO(\mu^3e^{-2\mu/3}).
\end{equation}
In particular the smooth-bias residual is eventually negative on even
indices and positive on odd indices.
\end{corollary}
\begin{proof}
Take $K=2$, use $A_2(n)=(-1)^n$, and substitute
$\nu=\mu(1-z^2v)$. This simplifies the prefactor in
\eqref{eq:sectors} to the amplitude in \eqref{eq:amplitude}.
The quotient of the remainder by the leading exponential is
$\OO(\mu^2e^{-\mu/6})\to0$, while $\mathcal A(1/\mu)\to1$.
\end{proof}

The first arithmetic layer explains why a smooth large-$n$ expansion
can miss small odd-index behavior. It does not by itself prove the
finite exceptions, which were separately certified.

There is an important limitation to the word ``transseries'' here.
The actions $1-1/k$ in \eqref{eq:sectors} accumulate at $1$, and
nonlinear products of the first exponential layer also reach action
$1$. Theorem~\ref{thm:sectors} is a rigorous \emph{fixed finite-sector}
statement below that accumulation point. It is not a proof that a
naively ordered infinite formal transseries remains valid across the
accumulation. Developing an appropriate resummation there is a genuine
further question.

\section{Consequences for differences and perturbations}\label{sec:consequences}
\subsection{Fixed-order eventual alternating differences}
Write $\Delta f(n)=f(n+1)-f(n)$ and let
$(1/2)_r=(1/2)(3/2)\cdots(r-1/2)$, with $(1/2)_0=1$.

\begin{theorem}[Every fixed finite-difference order]\label{thm:differences}
For every fixed integer $r\ge0$,
\begin{equation}\label{eq:differenceasym}
 (-1)^r\Delta^r(d_n-\delta)
 \sim\frac5{c^3}(1/2)_r\,n^{-r-1/2}
 \qquad(n\to\infty).
\end{equation}
Consequently $d_n$ is eventually strictly decreasing and strictly
convex, and every fixed-order alternating difference is eventually
positive.
\end{theorem}
\begin{proof}
Let $f(x)=B((c\sqrt{x-a})^{-1})-\delta$. The convergent analytic
series near infinity gives, for fixed $r$,
\[
 f^{(r)}(x)=(-1)^r\frac5{c^3}(1/2)_r x^{-r-1/2}
 +\OO_r(x^{-r-1}).
\]
Termwise differentiation is justified by local uniform convergence on
a complex neighborhood of each sufficiently large positive ray. For
$r\ge1$,
\[
 \Delta^r f(n)=\int_{[0,1]^r}
 f^{(r)}(n+t_1+\cdots+t_r)\,\dd t_1\cdots\dd t_r.
\]
The case $r=0$ is the original expansion. On the other hand,
Theorem~\ref{thm:effective} bounds the discrete remainder
$d_n-\delta-f(n)$ by an exponentially small function of $\sqrt n$.
Its $r$th finite difference is bounded by the sum of $2^r$ such local
bounds, which is smaller than every fixed negative power of $n$.
Combining these facts proves \eqref{eq:differenceasym}.
\end{proof}

The order of quantifiers matters: for each fixed $r$ there exists a
threshold $N_r$. No single threshold valid for all $r$ is claimed.
The small-index increase from $d_{10}$ to $d_{11}$ also shows why
``eventually decreasing'' cannot simply be strengthened to monotonicity
throughout the OEIS conjecture's range.

\subsection{The natural noise scale for recovering an index}
\begin{proposition}[Inverse sensitivity at the critical scale]\label{prop:noise}
Let $\eta\in\R$ be fixed and suppose
\begin{equation}\label{eq:noisyinput}
 y_n=p(n)\exp\!\left(\frac\eta{\sqrt n}+o(n^{-1/2})\right).
\end{equation}
Then
\begin{equation}\label{eq:noiselimit}
 F(y_n)-n\longrightarrow-\delta+\frac{2\eta}{c}.
\end{equation}
In particular, nearest-integer recovery eventually succeeds if
\begin{equation}\label{eq:noiserange}
 \frac c2(\delta-\tfrac12)<\eta<
 \frac c2(\delta+\tfrac12).
\end{equation}
It eventually fails if the limit in \eqref{eq:noiselimit} lies outside
$[-1/2,1/2]$. For the bias-corrected inverse $F(y_n)+\delta$, the
corresponding sufficient condition is $|\eta|<c/4$.
\end{proposition}
\begin{proof}
On the logarithmic target scale, the derivative in
\eqref{eq:PhiDerivative}, now evaluated near $\log p(n)$, is
$(2/c)\sqrt n+\OO(1)$. The mean value theorem applied to
\eqref{eq:noisyinput} therefore gives
$F(y_n)-F(p(n))\to2\eta/c$. Subtracting $d_n\to\delta$ proves
\eqref{eq:noiselimit}. A limit strictly inside the rounding interval
implies eventual success, and a limit strictly outside its closure
implies eventual failure.
\end{proof}

At either boundary of \eqref{eq:noiserange}, the leading limit alone
does not decide success; the next error term and the tie convention
can matter. The proposition deliberately leaves these boundary cases
unasserted.

\subsection{Relation to the forward approximation in A050811}
The sequence A050811 rounds $H(n)$ itself and records an asymptotic
formula for its difference from $p(n)$ \cite{oeis050811}. Our inverse
bias is consistent with that forward error. Indeed,
$F(p(n))=n-d_n$ and $d_n\to\delta$ imply
\[
 \log\frac{H(n)}{p(n)}
 =\log H(n)-\log H(n-d_n)
 \sim\frac{c\delta}{2\sqrt n}.
\]
Consequently
\begin{equation}\label{eq:forwarderror}
 H(n)-p(n)\sim
 \left(\frac1\pi+\frac\pi{72}\right)
 \frac{e^{\pi\sqrt{2n/3}}}{4\sqrt2\,n^{3/2}}.
\end{equation}
Rounding $H(n)$ changes its value by at most $1/2$, negligible on this
scale, so the same leading equivalent applies to A050811$(n)-p(n)$.
The leading equivalent in \eqref{eq:forwarderror} is already present
in that OEIS entry; it is included as a consistency check, not as a
novel formula. In particular, a small and bounded inverse error is
compatible with an unbounded forward error.

\section{Reproducibility and numerical illustrations}\label{sec:computations}
\subsection{Three deliberately separate forms of verification}
The archive separates exact proof checks from illustrative numerical
calculations.

\paragraph{Finite proof certificate.}
\texttt{code/certify.py} uses only the Python standard library, with
integer and rational arithmetic. It verifies
Lemma~\ref{lem:finite}, independently recomputes 400 partition values,
and encloses the sharp constants. Its output is
\texttt{results/certificate.json}; the companion CSV lists all 1,204
interval comparisons with exact endpoints. The analytic large-index proof is
not encoded in this program.

\paragraph{Exact coefficient calculation.}
\texttt{code/series.py} evaluates the triangular recurrences over the
rational numbers. The supplied run computes order 12 and writes
\texttt{results/coefficients.json}. Both defining logarithmic equations are also verified by exact
substitution through degree 13. Its output is exact algebra,
not a high-precision fit to partition data.

\paragraph{Numerical experiments.}
\texttt{code/experiments.py} uses \texttt{mpmath} with 130 decimal
digits. It evaluates the Lambert inverse, the analytic fixed points,
the parity residual and the staircase formula. The resulting
\texttt{results/numerics.json} is explicitly labeled numerical, not
an interval certificate. These experiments check implementation and
illustrate scale; they are not substitutes for the analytic proofs.

From the archive root, the reproduction commands are
\begin{lstlisting}
python code/certify.py --output results/certificate.json
python code/series.py --order 12 --output results/coefficients.json
python code/experiments.py --output results/numerics.json
pdflatex -halt-on-error article.tex
pdflatex -halt-on-error article.tex
\end{lstlisting}
Only the third command needs the optional \texttt{mpmath} dependency.
The \texttt{Makefile} automates these targets. The certificate must be
run without Python's assertion-removing optimization flag.

\subsection{The small indices that decide the theorem}
Table~\ref{tab:bias} gives rounded decimal illustrations. Their role is
expository; all strict comparisons used in the proof were separately
certified with rational intervals.
\begin{center}
\begin{longtable}{@{}r r l c@{}}
\caption{Inverse bias and nearest-integer recovery.}\label{tab:bias}\\
\toprule
$n$ & $p(n)$ & $d_n$ (approximately) & Recovery\\
\midrule
\endfirsthead
\toprule
$n$ & $p(n)$ & $d_n$ (approximately) & Recovery\\
\midrule
\endhead
2 & 2 & 0.796615481265 & fails\\
3 & 3 & 0.756104945769 & fails\\
4 & 5 & 0.502652583555 & fails\\
5 & 7 & 0.645462483764 & fails\\
6 & 11 & 0.438137374848 & succeeds\\
7 & 15 & 0.569463332725 & fails\\
8 & 22 & 0.450004070600 & succeeds\\
9 & 30 & 0.505033636323 & fails\\
10 & 42 & 0.440766683161 & succeeds\\
11 & 56 & 0.498487575397 & succeeds\\
12 & 77 & 0.420669769623 & succeeds\\
13 & 101 & 0.473632529547 & succeeds\\
14 & 135 & 0.431618414918 & succeeds\\
15 & 176 & 0.452838038452 & succeeds\\
\bottomrule
\end{longtable}
\end{center}
The values at $4$ and $9$ lie just above $1/2$; the value at $11$ lies
just below. An argument based only on a few leading terms at those
indices would be unreliable.

\subsection{Checking the parity term and the convergent tails}
Let
\[
 R_n=d_n-B(1/\mu),\qquad
 L_n=-(-1)^n\frac{\sqrt2\mu}{c^2}
                  \mathcal A(1/\mu)e^{-\mu/2}.
\]
The quotient $R_n/L_n$ should tend to $1$ by
Corollary~\ref{cor:parity}. The following rounded values use the full
analytic fixed point in $B$ and $\mathcal A$, not a truncated polynomial.
\begin{center}
\begin{tabular}{@{}r r r@{}}
\toprule
$n$ & $R_n$ (approximately) & $R_n/L_n$ (approximately)\\
\midrule
100 & $-1.5380713277\cdot10^{-5}$ & $0.995255662344$\\
400 & $-8.1517489953\cdot10^{-11}$ & $0.999897873292$\\
401 & $+7.9066922485\cdot10^{-11}$ & $1.000195579199$\\
1000 & $-4.2960209400\cdot10^{-17}$ & $0.999999260234$\\
1001 & $+4.2118871198\cdot10^{-17}$ & $1.000001383952$\\
2000 & $-3.0616805313\cdot10^{-24}$ & $0.999999994825$\\
2001 & $+3.0188441869\cdot10^{-24}$ & $0.999999992109$\\
\bottomrule
\end{tabular}
\end{center}
The oscillation of the ratios about $1$ reflects the next arithmetic
sectors, not a failure of the leading parity prediction.

For the twelve-term analytic bias approximation, the observed absolute
errors in $B(1/\mu)$ at $n=400,1000,2000$ are approximately
$4.81\cdot10^{-18}$, $1.22\cdot10^{-20}$ and $1.33\cdot10^{-22}$.
These are much smaller than the deliberately conservative Cauchy bounds.
At $n=2000$, the algebraic truncation error after twelve terms is larger
than the leading parity residual. This illustrates why an exponentially
small signal cannot automatically be read off after a fixed short
power-series truncation, even at high arithmetic precision.

For the corrected inverse, the experiments evaluate the twelve-term
version of \eqref{eq:Gexact} at the target $M(n)$, where the exact answer
is $n$. The errors (exact value minus approximation) at those same indices are about
$-2.46\cdot10^{-19}$, $-6.22\cdot10^{-22}$ and
$-6.79\cdot10^{-24}$. This tests the \emph{inverse} coefficient family
independently of the forward-bias coefficients.

\section{Further research questions and topics}\label{sec:research}
The following questions are proposed continuations of the proved
results. Except for the OEIS conjecture resolved above, no claim is
made that every formulation has already appeared as an open problem
in the literature.

\paragraph{1. Locate the nearest complex singularities.}
The radius $1/20$ is a convenient sufficient radius, not an optimal one.
Determine the exact radii of convergence of $v$ and $q$, the dominant
singularities on their convergence circles, and the resulting large-order
asymptotics of $w_j$ and $r_j$. The fixed-point equations suggest studying
simultaneous vanishing of the implicit equation and its derivative in
the unknown function. A rigorous answer would give much sharper
coefficient and truncation bounds than the elementary disk estimates.

\paragraph{2. Prove or refute positivity of every bias coefficient.}
All $v_j$ are nonnegative by the triangular recurrence, and the first
twelve $w_j$ are positive. Determine whether $w_j>0$ for every $j\ge1$.
If true, seek a positive recurrence or a combinatorial interpretation
that explains the cancellation in $2v-z^2v^2$. Such a result could
produce monotone lower approximants to the smooth bias.
The sign-changing corrected-inverse coefficients $r_j$ provide a useful
contrast: positivity must not be inferred merely from invertibility.

\paragraph{3. Find sharp thresholds for finite differences.}
For fixed $r$, Theorem~\ref{thm:differences} proves eventual positivity
of $(-1)^r\Delta^r(d_n-\delta)$. Determine the least threshold $N_r$,
starting with monotonicity and convexity, and study the growth of $N_r$
with $r$. The parity term suggests that a joint regime with increasing
$r$ is substantially different from the fixed-$r$ theorem. An effective
proof should track both differentiation of the analytic sector and the
amplification of alternating arithmetic errors.

\paragraph{4. Understand the extremum on every terminal range.}
For an integer $N\ge2$, define
$\beta(N)=\sup_{n\ge N}d_n$. We determined it sharply for $N=2$ and
$N=10$. Classify the maximizing index as a function of $N$, identify
where that index changes, and derive an effective description of the
optimal shift $(\beta(N)+\delta)/2$. Eventual monotonicity would imply
$\beta(N)=d_N$ beyond a threshold, but does not describe the finite
transition region.

\paragraph{5. Resolve the accumulation of exponential actions.}
The sector actions $1-1/k$ accumulate at $1$, while the square of the
parity correction has action $1$. Construct a mathematically suitable
resummation or grouped expansion through this level. The target should
be a theorem with a uniform remainder when the number of retained
Rademacher terms grows with $n$, not a formal rearrangement of infinitely
many exponentials. The exact convergent Rademacher series supplies a
natural starting point.

\paragraph{6. Sharpen inverse algorithms and their certification.}
Replace the simple two-candidate theorem by an adaptive method that
combines the convergent correction series, interval Lambert evaluation
and an effective Rademacher routine. Prove bit-complexity bounds in the
size of the target, including inputs near the floor boundaries in
\eqref{eq:k}. The one-comparison result alone does not resolve the
cost of certifying that floor. Another concrete task is to determine
the sharp threshold for rounding the corrected model $G(p(n))$.

\paragraph{7. Extend exact inverse rounding beyond ordinary partitions.}
For partitions into distinct parts, colored partitions, or restricted
partition families with controlled exponential expansions, identify the
correct inverse scaling, limiting bias and possible successful constant
shifts. The appropriate analogue of the present theorem is a global
rounding statement with explicit exceptional indices, not only an
inverse asymptotic equivalent. Families whose forward relative error
falls exactly at the inverse sensitivity scale are especially natural.

\paragraph{8. Analyze the critical perturbation boundaries.}
At $\eta=c(\delta\pm1/2)/2$, Proposition~\ref{prop:noise} leaves a
half-integer limiting displacement. Add the next perturbation term,
for example $\theta/n$, and classify eventual success or failure.
After sufficiently many algebraic terms are canceled, the arithmetic
parity term can become the first decision-making term. This offers a
precise setting in which a beyond-all-orders correction controls an
integer-valued outcome.

\paragraph{9. Extract residue-sensitive error bounds.}
The $k=2$ term distinguishes even and odd indices. Incorporate the
$k=3,4,\ldots$ terms to obtain sharp inequalities on fixed residue
classes and to reduce the analytic-to-finite cutoff. A useful deliverable
would be a small rational certificate for the least possible cutoff
in an explicit proof architecture, rather than merely a longer table
of numerical confirmations.

\paragraph{10. Formalize the proof and expose its dependencies.}
Formalize the interval operations, the finite partition product and
Lemma~\ref{lem:finite} in Lean or Rocq. Separately formalize the analytic
contraction and inverse-error arguments. The large external dependency
is Rademacher's formula: a formal development should clearly distinguish
assuming an already-established statement of that theorem from
formalizing its proof. A verified certificate checker would reduce the
trusted code for the finite step without disguising that analytic
boundary.

\section{Conclusion}
The conjecture in A306631 is an exact-recovery phenomenon that can be
proved from classical partition analysis, but it is not a consequence
of the leading asymptotic equivalent alone. The bias is bounded away
from zero, approaches $1/24+3/\pi^2$, and has small-index arithmetic
fluctuations decisive for rounding. The combination of an explicit
infinite-range estimate and an exact rational finite certificate
resolves the conjecture and identifies its complete exception set.

The sharp extremum at $11$ then determines an optimal uniform shift,
a numerical safety margin, and a two-candidate staircase inversion.
The convergent analytic corrections and the separate arithmetic
exponential layers explain how these discrete facts sit inside a more
precise inverse theory. These are the concrete mathematical outcomes
of the report; neither a sweeping priority claim nor a general solution
to partition-function inversion is required for them.

\appendix
\section{A compact proof dependency ledger}\label{app:ledger}
\begin{center}
\begin{longtable}{@{}p{0.30\linewidth}p{0.62\linewidth}@{}}
\toprule
Result or object & Dependency and status\\
\midrule
\endfirsthead
\toprule
Result or object & Dependency and status\\
\midrule
\endhead
Rademacher identity & Quoted classical theorem, in the normalization of
Johansson \cite{johansson}. Not proved by the code.\\
Elementary tail & Derived here from $|A_k(n)|\le k$ and an integral
bound on $J(t)$.\\
Convergent bias model & Analytic contraction, exact algebra and Cauchy
estimates, proved here.\\
All $n\ge400$ & Relative Rademacher error, derivative lower bound and
an explicit bracketing argument, proved here.\\
$2\le n\le399$ & Exhaustive strict comparisons with rational interval
endpoints; standard-library certificate supplied and executed.\\
Partition data & Exact finite-product computation independently checks
all values used by the finite certificate.\\
Sharp rounding and extrema & Analytic infinite range plus the exact
finite lemma.\\
Optimal shifts and staircase & Consequences of the sharp bias extrema
and monotonicity.\\
Inverse coefficient lists & Exact rational recurrences, not numerical
fitting.\\
Arithmetic corrections & Exact Rademacher sectors and Taylor expansion
of the logarithmic inverse; fixed finite number of sectors.\\
Numerical tables & High-precision illustrations; not used as proof of
strict inequalities.\\
Priority and formal status & Unrefereed report; limited search; no
Lean/Rocq kernel checking claimed.\\
\bottomrule
\end{longtable}
\end{center}

\section{Suggested OEIS update and source audit}\label{app:oeis}
An editor-ready draft is included as \texttt{OEIS\_update\_draft.txt}.
Its central replacement for the conjectural statement is the following
mathematical content, subject to independent editorial review:
\begin{quote}
For every integer $n\ge10$, A306631$(p(n))=n$. For $n\ge2$, the only
exceptions to this identity are $2,3,4,5,7,9$, and at each exception the
value is $n-1$.
\end{quote}
A second useful comment would state the ceiling identity and the sharp
bias bounds, with a reference to a stable, publicly available version
of this report after review. No OEIS submission was made as part of
this work.

The source audit distinguished the entry's own revision marker from
the site-wide footer date. The A306631 internal page inspected for this
report displayed its formula and conjecture, with entry revision
\texttt{\#16, Feb 16 2025}. The report date records when that content was
checked, not a claim that the conjecture was first posted on that date.
The entry credits Jean-Fran\c{c}ois Alcover and dates its origin to
2 March 2019 \cite{oeis306631}.

The limited bibliographic check did not locate a prior proof of this
specific exact-rounding statement. That negative search result is not
a proof of novelty. The mathematical assertions are therefore stated
as theorems proved in this report, with classical inputs credited,
without claiming that every reformulation or consequence is
historically unprecedented.

\begin{thebibliography}{9}
\bibitem{oeis306631}
OEIS Foundation Inc., \emph{A306631: Inverse of the Hardy--Ramanujan
asymptotic partition function}. Entry credited to Jean-Fran\c{c}ois
Alcover, 2 March 2019. Internal revision \#16 displayed 16 February 2025;
consulted 3 October 2026.
\url{https://oeis.org/A306631/internal}.

\bibitem{oeis000041}
OEIS Foundation Inc., \emph{A000041: Number of partitions of $n$}.
\url{https://oeis.org/A000041}.

\bibitem{oeis050811}
OEIS Foundation Inc., \emph{A050811: Partition numbers rounded to nearest
integer given by the Hardy--Ramanujan approximate formula}.
Consulted 3 October 2026.
\url{https://oeis.org/A050811/internal}.

\bibitem{johansson}
F. Johansson, \emph{Efficient implementation of the
Hardy--Ramanujan--Rademacher formula}, LMS Journal of Computation and
Mathematics \textbf{15} (2012), 341--359.
\url{https://doi.org/10.1112/S1461157012001088}.
Preprint version 2:
\url{https://arxiv.org/abs/1205.5991v2}.

\bibitem{banerjee}
K. Banerjee, P. Paule, C.-S. Radu and C. Schneider,
\emph{Error bounds for the asymptotic expansion of the partition
function}, 2022, arXiv:2209.07887.
\url{https://arxiv.org/abs/2209.07887}.
Cited for context on forward error bounds, not as a hidden proof
assumption for the finite cutoff used here.

\bibitem{repo292507}
ProveIt repository maintained by V. Reshetnikov, \emph{All orders asymptotics and
asymptotic inversion for A292507}, report dated 1 October 2026, with
repository integration notes dated 2 October 2026. Inspected at commit
\texttt{d68b65ea064084fb121df9bb431b21d086f7be81}, path
\nolinkurl{SetTheory/Cardinals/docs/reports/generating-functions-and-asymptotics/oeis-sequence-asymptotics/a292507-binomial-partition-transform/article.tex}.
Repository: \url{https://github.com/VladimirReshetnikov/ProveIt}.
This repository source is methodological context, not a proof dependency.
\end{thebibliography}
\end{document}
